\section{Related work}
\label{sec:related}

Our study sits at the intersection of four literatures: ligand-based activity
prediction and the transfer assumption behind its evaluation, machine learning
for transition-metal complexes, few-shot and meta-learning, and evaluation
protocols for distribution shift. We take each in turn and close with the
specific gap we address.

\subsection{Ligand-based activity prediction and its transfer assumption}

Modern ligand-based activity prediction is dominated by message-passing graph
neural networks \citep{gilmer2017mpnn,hamilton2017graph}. The directed
message-passing neural network (D-MPNN) of \citet{dmpnn2019} remains a standard
backbone, and its reference implementation \citep{chemprop} is the de-facto
baseline on MoleculeNet \citep{moleculenet2018}. The family has since been
extended along three axes that matter here: attention over atom environments
\citep{attentivefp}, contrastive self-supervision \citep{molclr}, and
large-scale pretraining on chemical language \citep{molformer}. These sit inside
a broader programme of machine learning for chemistry
\citep{butler2018,mater2019} and molecular design
\citep{sanchezlengeling2018,gomezbombarelli2018}, whose activity-prediction
applications now extend to prospectively validated discoveries
\citep{stokes2020}.

What all of these share is an evaluation convention: a \emph{scaffold split},
in which the Bemis--Murcko scaffold \citep{bemis1996murcko} of every test
compound is absent from training, so that the reported AUC is a transfer
estimate over organic chemical space. The convention encodes an assumption ---
that scaffold novelty is the hard case. Our question is orthogonal. We hold the
evaluation machinery fixed and vary a different axis: whether the \emph{metal
centre} of a coordination compound was observed during training. A scaffold
split of metal complexes does not control for this, because the same organic
scaffold can appear on several metals.

\subsection{Machine learning for transition-metal complexes}

Transition-metal complexes break several assumptions that organic benchmarks
rely on. Molecular graphs are ambiguous about bond order at the metal, valence
is variable, and the electronic structure that drives activity is not fully
captured by a 2D graph; the machine-learning approaches developed for this
chemical space are surveyed by \citet{nandy2021computational}. Concretely,
feature-selection studies for first-row transition-metal complexes
\citep{janet2017} and active-learning searches over spin-crossover space
\citep{janet2018} showed that quantum-chemical targets for coordination
compounds are learnable from structure alone.

Two strands of that work bear directly on our design. The first is pretraining
on quantum-chemistry corpora. The tmQM corpus \citep{tmqm2020} supplies quantum
geometries and properties for $\sim$86k transition-metal complexes, and
pretraining on it is a natural analogue of the large-scale pretraining that
lifted organic property prediction. It is one of a family of datasets that made
such pretraining credible --- from QM9 \citep{qm9} and ANI-1 \citep{ani1} to the
Open Catalyst Project \citep{opencatalyst2021} and crystal-graph models
\citep{cgcnn2018} --- and the tmQMg extension adds excited-state properties for
a further $74$k complexes \citep{tmqmg2025}. The second is metal-aware
descriptors. \citet{electrum2025} build a universal metal fingerprint from
electron configuration, which is exactly the kind of inductive bias a shared
ligand-only head lacks.

The metallodrug context that motivates the activity endpoint begins with
cisplatin \citep{rosenberg1969} and its cellular pharmacology
\citep{lippard2005}, and has since broadened to ruthenium \citep{alessio2007}
and organometallic \citep{zhangsadler2017} anticancer complexes; the drug-design
perspective is surveyed by \citet{ndagi2017} and the medicinal-inorganic framing
by \citet{mjos2014metallodrugs}.

On the activity side, \citet{krasnov2026metalcytotoxdb} assembled
MetalCytoToxDB, a curated database of anticancer transition-metal complexes
spanning Ru, Ir, Rh, Re and Os, and reported per-metal binary activity models
(Ru 0.81, Ir 0.73 AUC) trained on Morgan fingerprints with gradient-boosted
trees under a DOI-grouped 10-fold split. That work establishes both the data and
the protocol we adopt, and it is the baseline against which we measure our
within-metal comparison in \S\ref{sec:results:krasnov}.

\subsection{Few-shot and meta-learning}

Adapting a model to a new task from a handful of labelled examples is the
setting of few-shot learning, and its dominant formalisation is
gradient-based meta-learning: MAML \citep{maml2017} learns an initialisation
from which a few gradient steps adapt to a new task, and Reptile
\citep{reptile2018} simplifies that outer loop to a first-order update. In
molecular property prediction the idea has been applied as one-shot learning
over a shared embedding across many tasks \citep{oneshotdrug2017}, where the
scarcity is labels per task. Our setting differs in what plays the role of the
task. The tasks are \emph{metals}, the tasks share an organic ligand space, and
the shift between them is not a new label function over the same compounds but a
different conditional relationship between the ligand features and the label.
\fsma{} is deliberately weaker than a learned meta-learner --- a frozen source
encoder plus a retrained head, with no outer loop --- and we present it as a
strong baseline for this axis rather than a new meta-learning algorithm. The
result that matters is that even this minimal adaptation is enough to repair the
collapse, which is the evidence that the collapse is a target-distribution
problem rather than a representation one.

\subsection{Evaluation protocols and distribution shift}
\label{sec:related:shift}

Our finding is a statement about evaluation protocol, so it also belongs to the
literature on how molecular models are benchmarked and where they fail. Domain
shift in molecular machine learning is usually studied as a shift in scaffold,
size, or assay, and is benchmarked explicitly by datasets such as DrugOOD
\citep{drugood2023}; the general problem is studied as domain generalisation,
for which comprehensive surveys exist \citep{zhoudg2023,wangdg2021}, and the
closely related covariate-shift setting has its own estimators and validation
schemes \citep{sugiyama2006}. The broader
graph-learning-for-drug-discovery toolkit is surveyed by \citet{gnn2025review},
and the question of when a deep QSAR model can be trusted outside its training
distribution --- its applicability domain --- is examined by
\citet{applicabilitydomain2019}.

The shift we study is not of that kind. It is a change in the identity of a
single atom whose chemistry reorganises the whole molecule, so the test
distribution is not a novel point in the same chemical space but a different
conditional relationship between the ligand features and the label. We are not
aware of an organic benchmark that isolates this effect, which is precisely why
we treat it as a finding rather than an implementation detail.

\subsection{What is missing}
\label{sec:related:gap}

Across both the ligand-based and the transition-metal literature, transfer is
reported \emph{within} a metal or across organic chemical space. A per-metal
model answers ``how well can we predict activity for this metal''; it does not
answer ``does what we learn about one metal help with another''. No study in
either line publishes a cross-metal leave-one-metal-out measurement, and none
reports a repair that is validated against budget-matched controls and a
different model class. We make cross-metal transfer the primary measurement,
report a minimal adaptation protocol that recovers it, and establish that the
recovery is the adaptation rather than the extra labels, the extra compute, or
the choice of encoder.
